%%%%%%%%%%%%%%%%%%%%%%%%%%%%%%%%%%%%%%%%%%%%%%%%%%%%%%%%%%%%%%%%%%%%%%%%%%%
% Lehreinheit 6: Forschungsdesign
% Quelle: WuZ-Buch, Kap. 15 (LE 6), Kap. 12d (Videostudie Becher-Monty-Hall:
%         Studienlogik, Laengendifferenz), Kap. 12f (Feldstudie, Endbefunde),
%         Kap. 10b (Testbarkeit: Blind-Design, Prozessindikator,
%         Parallelmessung), Kap. 4 (HCEG als Validitaetsbedrohung)
% Staunmoment: Kontrollierte Seilroutine  |  Kernbegriff: interne Validitaet
% Geheimhaltung: Methode der Seilroutine wird NICHT beschrieben.
%%%%%%%%%%%%%%%%%%%%%%%%%%%%%%%%%%%%%%%%%%%%%%%%%%%%%%%%%%%%%%%%%%%%%%%%%%%
\nsl{In LE~5 wurde aus der Forschungsfrage eine Hypothese mit
Widerlegungsbedingung. Jetzt geht es um den Plan, mit dem man sie prüft.
Ein Forschungsdesign legt fest, wer was wann unter welchen Bedingungen
misst und womit verglichen wird. Die Einheit beginnt mit drei Seilen
unterschiedlicher Länge, die gleich lang werden. Die Gruppe sieht das
Kunststück zweimal: einmal so, wie es auf der Bühne gezeigt wird, und
einmal unter Bedingungen, die sie selbst festlegt. Der Unterschied zwischen
beiden Durchgängen ist der Unterschied zwischen einer Behauptung und einem
Experiment. Daraus entwickelt die Einheit Designtypen, interne und externe
Validität, Konfundierung und Kontrolle sowie die Gütekriterien von
Messungen. Als Fallbeispiel dient die Becher-Monty-Hall-Studie aus dem
Buch, mit einem sauberen Idealentwurf und einer Feldstudie, die zeigt, was
in der Praxis davon übrig bleibt.
\newpage}
\os{\begin{description}
\item[{\rot\bf Heute}]~\\[-\lblskip]
\bi
  \ite drei Seile, zwei Vorführungen
  \ite Forschungsdesign: Variablen, Vergleich, Kausalität
  \ite Designtypen: experimentell, quasi-experimentell, nicht-experimentell
  \ite interne und externe Validität, Konfundierung
  \ite Gütekriterien: Objektivität, Reliabilität, Validität
  \ite Fallbeispiel: die Becher-Monty-Hall-Studie
\ei
\end{description}
\newpage}

\begin{vorbereitung}[Kontrollierte Seilroutine]
\textbf{Eine Woche vorher}
\begin{checkliste}
  \item Seil-Set für eine Routine, in der drei Seile unterschiedlicher
    Länge gleich lang werden (klassischer Effekt, im Fachhandel als
    fertiges Set erhältlich oder nach der Standardliteratur selbst
    zusammengestellt). Zwei vollständige Sets einpacken; Seile werden
    schnell unansehnlich.
  \item \textbf{Methode und Einstudierung nach der Anleitung des
    Herstellers bzw.\ der Standardliteratur; hier bewusst nicht
    beschrieben.} Das Skript wird öffentlich gespiegelt.
  \item Zwei Präsentationen proben. \emph{Variante~A} wie auf der Bühne:
    Seile zeigen, Längen behaupten, Effekt, Seile ins Publikum. Dauer etwa
    drei Minuten. \emph{Variante~B} mit Kontrollbedingung: Die Gruppe legt
    vorab fest, wie Ausgangs- und Endlängen geprüft werden.
  \item Vor der Probe von Variante~B entscheiden, welche Prüfungen Sie
    zulassen und welche Sie ausdrücklich verweigern. Faustregel: Alles,
    was die Methode gefährden könnte, wird freundlich und offen abgelehnt
    (\enquote{Diese Kontrolle lasse ich nicht zu.}). Die Ablehnung ist kein
    Notbehelf, sondern der didaktische Kern der Variante.
  \item Für Variante~B ein Maßband oder einen Zollstock (2\,m) und einen
    Flipchart-Bogen für das Messprotokoll bereitlegen.
  \item Das Fallbeispiel vorbereiten: \buch{Kap.~12d} (Studienlogik,
    Beat-Plan, Längendifferenz) und \buch{Kap.~12f} (Feldstudie,
    Ergebnisse). Zahlen nur so verwenden, wie sie dort stehen.
\end{checkliste}

\textbf{Handout und Übungsblatt}
\begin{checkliste}
  \item Handout \enquote{Forschungsdesign} und Übungsblatt LE~6 je einmal
    pro Person kopieren (Druckvorlage \texttt{handouts.pdf}, Seiten zu
    LE~6). Das Übungsblatt enthält die fiktive Studienbeschreibung für die
    Design-Analyse und die Vorlage für die Design-Skizze vollständig.
  \item Die H\textsubscript{0}/H\textsubscript{1}-Blätter aus LE~5
    (Aufgabe~3) zurückgeben bzw.\ mitbringen lassen; sie sind die
    Grundlage der eigenen Design-Skizze.
  \item Optional statt des fiktiven Falls einen echten, kurzen
    empirischen Artikel aus dem Fach der Gruppe auswählen (Methodenteil
    höchstens zwei Seiten, klare Gruppenvergleiche). Dann den
    Methodenteil pro Zweiergruppe einmal kopieren.
\end{checkliste}

\textbf{Am Tag selbst}
\begin{checkliste}
  \item Seile vor der Stunde vorbereiten; das Ersatzset griffbereit, aber
    außer Sicht.
  \item Tafel in zwei Spalten teilen: \enquote{Variante A: Was haben wir
    gesehen?} und \enquote{Variante B: Was haben wir gemessen?}.
  \item Für die Design-Werkstatt Zweiergruppen einteilen, die sich in der
    Feedback-Phase gegenseitig die Skizzen vorstellen.
\end{checkliste}

\textbf{Risiken und Plan B.} Die Seilroutine ist bekannt; es kann sein,
dass jemand sie schon gesehen hat oder die Methode kennt. Diese Person
bitten, ihr Wissen für sich zu behalten, und den Umstand selbst nutzen:
Vorwissen ist eine Störvariable, und genau sie hat in der
Monty-Hall-Feldstudie alle Armunterschiede überstrahlt (siehe
Fallbeispiel). Misslingt der Effekt, wird er nicht wiederholt; dann
trägt Variante~B die Stunde allein, indem die Gruppe ein Messprotokoll für
eine Behauptung entwirft, die die Lehrperson nur ankündigt. Ohne Seile
lässt sich der Einstieg mit einer bekannten Längentäuschung
(Müller-Lyer-Figur) erst ohne, dann mit Lineal ersetzen.
\end{vorbereitung}

\begin{ablauf}
\begin{description}
\item[0--5 · Blitzlicht.] Von der Hypothese (LE~5) zum Prüfplan:
  \enquote{Wie würden Sie Ihre Hypothese testen?} Zwei, drei Antworten
  sammeln.
\item[5--15 · Staunmoment.] Seilroutine erst als Variante~A ohne
  Vergleich, dann als Variante~B mit Kontrollbedingung, die die Gruppe
  festlegt. Die Gruppe erlebt, wie die Kontrolle die Täuschung entlarvt
  oder ihre Verweigerung die Behauptung entwertet.
\item[15--45 · Theorie-Input.] Designtypen; interne und externe
  Validität; Konfundierung; Gütekriterien. Fallbeispiel: das Design der
  Becher-Monty-Hall-Videostudie und ihr Störfaktor, dann die Feldstudie.
\item[45--75 · Design-Werkstatt.] Analyse der fiktiven Studie (UV, AV,
  Kontrollgruppe, Konfounder; Übungsblatt, Aufgabe~1); dann Skizze eines
  eigenen Designs zur eigenen Hypothese (Aufgabe~2); Messgüte-Check mit
  dem Vokabular aus \buch{Kap.~10b}.
\item[75--90 · Sicherung.] Peer-Feedback auf die Design-Skizzen in
  Zweiergruppen; Abgabe der Skizze als benotete Teilleistung
  terminieren.
\end{description}
\end{ablauf}

%==============================================================================
\section{Staunmoment: Drei Seile}
%==============================================================================
\begin{staunbox}{Kontrollierte Seilroutine}
Drei Seile: kurz, mittel, lang. \emph{Ich mache sie gleich lang.}

\medskip
\textbf{Variante A:} Schauen Sie zu.

\textbf{Variante B:} Diesmal bestimmen Sie, wie wir prüfen. Was wollen Sie
vorher messen, was nachher, und womit vergleichen?
\end{staunbox}
\os{\newpage}

\begin{description}
\item[{\rot\bf Variante A: ohne Vergleich}]~\\[-\lblskip]
\bi
  \ite Was haben Sie gesehen?
  \ite Woher wissen Sie, dass die Seile vorher verschieden lang waren?
  \ite Woher wissen Sie, dass sie nachher gleich lang sind?
  \ite Wer hat gemessen -- und womit?
\ei
\os{\newpage}

\item[{\rot\bf Variante B: mit Kontrollbedingung}]~\\[-\lblskip]
\bi
  \ite Ihr Prüfplan: Was messen wir vorher, was nachher?
  \ite Wer misst? Wer darf die Seile halten?
  \ite Womit vergleichen wir: Maßband, ein unbeteiligtes Seil, eine
    zweite Person?
  \ite Welche Prüfung lässt die Vorführende \emph{nicht} zu -- und was
    folgt daraus?
\ei
\os{\newpage}
\end{description}

\nsl{In Variante~A ist die einzige Evidenz die Vorführung selbst. Die
Längen vorher und nachher werden behauptet und durch Zeigen
\enquote{belegt}; gemessen hat niemand, verglichen wurde nichts. In
Variante~B formuliert die Gruppe einen Prüfplan, bevor das Kunststück
beginnt. Lässt die Lehrperson eine Prüfung zu, wird der Effekt an einem
unabhängigen Maßstab gemessen. Verweigert sie eine Prüfung, wird sichtbar,
dass die Behauptung nur unter Bedingungen gilt, die der Behauptende selbst
kontrolliert. Für ein Kunststück ist das legitim, denn es behauptet
nichts über die Welt. Für eine Studie wäre es ein K.-o.-Kriterium.
Dieselbe Figur kennt das Buch vom Muskeltest der Applied Kinesiology:
Solange der Anwender den Test kontrolliert, scheint er zu funktionieren;
in einer Doppelblindstudie lag die Trefferquote kaum über dem Zufall
(\buch{Kap.~10b}).}

\begin{description}
\item[{\rot\bf Was ist der Unterschied?}]~\\[-\lblskip]
\bi
  \ite A: Behauptung plus Vorführung -- der Behauptende kontrolliert alles
  \ite B: Prüfplan vorab, unabhängiger Maßstab, Vergleich vorher/nachher
  \ite eine verweigerte Kontrolle ist selbst ein Befund
  \ite ohne Kontrollbedingung bleiben selbst grobe Täuschungen
    unentdeckt
\ei
\os{\newpage}
\end{description}

%==============================================================================
\section{Was ein Forschungsdesign leistet}
%==============================================================================
\nsl{Ein Forschungsdesign ist der Plan, mit dem eine Forschungsfrage
beantwortet werden soll. Es legt fest, welche Variablen erhoben oder
verändert werden, an wem, zu welchen Zeitpunkten und unter welchen
Bedingungen, und vor allem, womit verglichen wird
\parencite{doering2016forschungsmethoden}. Die \emph{unabhängige Variable}
(UV) ist die vermutete Ursache; im Experiment wird sie gezielt variiert.
Die \emph{abhängige Variable} (AV) ist die gemessene Wirkung.
\emph{Störvariablen} beeinflussen die AV ebenfalls, ohne Gegenstand der
Untersuchung zu sein. Will man aus einem Zusammenhang auf Kausalität
schließen, müssen drei Bedingungen erfüllt sein: Ursache und Wirkung
variieren gemeinsam, die Ursache geht der Wirkung zeitlich voraus, und
plausible Alternativerklärungen sind ausgeschlossen
\parencite{shadish2002experimental}. Die dritte Bedingung ist die
schwierigste, und an ihr entscheidet sich die Qualität eines Designs.}
\begin{description}
\item[{\rot\bf Forschungsdesign}]~
\end{description}
\defn{design}{Ein \textbf{Forschungsdesign} ist der Plan einer
Untersuchung: welche Variablen an wem, wann und unter welchen Bedingungen
erhoben oder verändert werden und womit die Ergebnisse verglichen werden.}
\bi
  \ite \textbf{UV:} vermutete Ursache (im Experiment gezielt variiert)
  \ite \textbf{AV:} gemessene Wirkung
  \ite \textbf{Störvariable:} wirkt ebenfalls auf die AV
\ei
\os{\newpage}

\begin{description}
\item[{\rot\bf Wann darf man auf Kausalität schließen?}]~\\[-\lblskip]
\bi
  \ite \textbf{Kovariation:} Ursache und Wirkung treten gemeinsam auf
  \ite \textbf{zeitliche Folge:} die Ursache geht der Wirkung voraus
  \ite \textbf{keine Alternativerklärung:} nichts anderes erklärt den
    Zusammenhang ebenso gut
  \ite die dritte Bedingung ist die schwierigste -- sie ist der Zweck des
    Designs
\ei
\os{\newpage}

\item[{\rot\bf Seile als Design}]~\\[-\lblskip]
\bi
  \ite UV: das Kunststück (durchgeführt / nicht durchgeführt)
  \ite AV: die Seillängen
  \ite Variante~A: eine Messung, nachher, durch den Behauptenden
  \ite Variante~B: Vorher- und Nachher-Messung, unabhängiger Maßstab
  \ite Welche Alternativerklärungen schließt B aus, A nicht?
\ei
\os{\newpage}
\end{description}

%==============================================================================
\section{Designtypen}
%==============================================================================
\nsl{Designs lassen sich danach ordnen, ob die UV aktiv variiert wird und
ob die Zuteilung zu den Bedingungen zufällig erfolgt
\parencite{shadish2002experimental, doering2016forschungsmethoden}. Im
\emph{Experiment} variieren die Forschenden die UV und teilen die
Teilnehmenden zufällig den Bedingungen zu (Randomisierung). Zufall
verteilt bekannte und unbekannte Störvariablen im Mittel gleich auf die
Gruppen; das macht das Experiment zum stärksten Design für
Kausalaussagen. Im \emph{Quasi-Experiment} wird die UV variiert, aber
vorhandene Gruppen werden verglichen (zwei Kurse, zwei Standorte); die
Gruppen können sich schon vorher unterschieden haben. Im
\emph{nicht-experimentellen} Design wird nichts verändert, nur erhoben;
es zeigt Zusammenhänge, aber kaum Ursachen. Quer dazu liegt die Frage nach
der Zeit: Ein \emph{Querschnitt} misst einmal und vergleicht Personen, ein
\emph{Längsschnitt} misst dieselben Personen mehrfach und kann
Veränderungen verfolgen. Die Notation von Campbell kürzt Designs ab:
R steht für zufällige Zuteilung, X für die Intervention, O für eine
Messung (\emph{observation}).}
\begin{description}
\item[{\rot\bf Drei Grundtypen}]~\\[-\lblskip]
\begin{center}
\renewcommand{\arraystretch}{1.25}
\begin{tabular}{p{0.24\textwidth}p{0.17\textwidth}p{0.17\textwidth}p{0.26\textwidth}}
\toprule
& \textbf{UV variiert?} & \textbf{Zufall?} & \textbf{Kausalaussage} \\
\midrule
Experiment & ja & ja & stark \\
Quasi-Experiment & ja & nein & eingeschränkt \\
nicht-experimentell & nein & -- & schwach (Zusammenhang) \\
\bottomrule
\end{tabular}
\end{center}
\bi
  \ite quer dazu: \textbf{Querschnitt} (einmal messen) oder
    \textbf{Längsschnitt} (dieselben Personen mehrfach)
\ei
\os{\newpage}

\item[{\rot\bf Designs in Kurzschrift}]~\\[-\lblskip]
\begin{center}
\renewcommand{\arraystretch}{1.3}
\begin{tabular}{p{0.46\textwidth}cccc}
\toprule
\textbf{Design} & & & & \textbf{Zeit $\rightarrow$} \\
\midrule
One-Shot (Variante A) & & & X & O \\
Vorher-Nachher, eine Gruppe (Variante B) & & O & X & O \\
\midrule
Experiment mit Kontrollgruppe & {\color{rwuviolet}R} & & X & O \\
 & {\color{rwuviolet}R} & & & O \\
\midrule
Experiment mit Vortest & {\color{rwuviolet}R} & O & X & O \\
 & {\color{rwuviolet}R} & O & & O \\
\bottomrule
\end{tabular}
\end{center}
\bi
  \ite R = Zufallszuteilung, X = Intervention, O = Messung
\ei
\os{\newpage}
\end{description}

%==============================================================================
\section{Interne Validität}
%==============================================================================
\nsl{Ein Design ist \emph{intern valide}, wenn sich die Veränderung der
AV eindeutig auf die UV zurückführen lässt. Shadish, Cook und Campbell
haben die typischen Bedrohungen zusammengestellt
\parencite{shadish2002experimental}. Die Liste ist ein Werkzeug: Wer ein
Design plant oder liest, geht sie Punkt für Punkt durch und fragt, welche
Bedrohung das Design ausschließt und welche nicht. Ein Vorher-Nachher-Vergleich
ohne Kontrollgruppe ist anfällig für fast alle; die zufällig zugeteilte
Kontrollgruppe neutralisiert die meisten, weil Geschichte, Reifung und
Testeffekte beide Gruppen gleich treffen.}
\begin{description}
\item[{\rot\bf Bedrohungen der internen Validität}]~\\[-\lblskip]
\begin{center}
\small
\renewcommand{\arraystretch}{1.15}
\begin{tabular}{p{0.24\textwidth}p{0.62\textwidth}}
\toprule
\textbf{Bedrohung} & \textbf{Beispiel: Lernvideo mit Vorher-Nachher-Test} \\
\midrule
Geschichte & in der Woche dazwischen läuft eine Doku zum Thema \\
Reifung & die Teilnehmenden lernen ohnehin dazu \\
Testeffekt & der Vortest übt schon für den Nachtest \\
Instrumentierung & Nachtest ist leichter als Vortest \\
Regression zur Mitte & nur die Schwächsten im Vortest werden ausgewählt \\
Selektion & Freiwillige mit Interesse gegen Pflichtgruppe \\
Ausfall & die Frustrierten brechen vor dem Nachtest ab \\
\bottomrule
\end{tabular}
\end{center}
\os{\newpage}
\end{description}

\nsl{Eine eigene Bedrohung entsteht durch die Forschenden selbst. Wenn
die Person, die das Ergebnis erhofft, auch misst oder bewertet, können
ihre Erwartungen die Daten verschieben, ohne dass sie es bemerkt. Der
\enquote{Kluge Hans}, ein Pferd, das angeblich rechnen konnte, antwortete
nur dann richtig, wenn die Fragenden die Antwort kannten und sichtbar
waren; es las ihre unwillkürlichen Bewegungen \parencite{pfungst1907}.
Dieselbe Logik steckt im ideomotorischen Pendel aus LE~1. Gegenmittel ist
das \emph{Blind-Design}: Wer misst, weiß nicht, in welcher Bedingung die
gemessene Person war. LE~7 vertieft diesen Versuchsleitereffekt.}
\begin{description}
\item[{\rot\bf Die Forschenden als Störvariable}]~\\[-\lblskip]
\bi
  \ite Der Kluge Hans (1907): ein Pferd \enquote{rechnet} -- solange die
    Fragenden die Antwort kennen und sichtbar sind
  \ite es reagierte auf unwillkürliche Bewegungen der Fragenden
  \ite Seile: wer vorführt, darf nicht auch messen
  \ite Gegenmittel: \textbf{Blind-Design} -- die Messenden kennen die
    Bedingung nicht
\ei
\os{\newpage}
\end{description}

%==============================================================================
\section{Konfundierung und Kontrolle}
%==============================================================================
\nsl{Von \emph{Konfundierung} spricht man, wenn eine Störvariable
systematisch mit der UV variiert. Dann lassen sich ihre Wirkungen nicht
mehr trennen. Wer eine Lern-App freiwillig nutzt, ist im Mittel vermutlich
motivierter als wer sie nicht nutzt; ein Notenunterschied kann von der App
oder von der Motivation kommen. Gegen Konfundierung gibt es fünf
Werkzeuge: Randomisierung, Konstanthalten (alle Bedingungen gleich bis auf
die UV), Parallelisieren oder Schichten (Gruppen nach einem bekannten
Störmerkmal ausbalancieren), eine geeignete Kontrollgruppe und
Blindbedingungen. Die Kontrollgruppe ist dabei nicht einfach \enquote{die
anderen}. Eine \emph{aktive Kontrolle} erhält eine gleich lange,
gleich aufwendige Bedingung ohne den vermuteten Wirkstoff, damit sich die
Gruppen nur in dem unterscheiden, was untersucht wird.}
\begin{description}
\item[{\rot\bf Konfundierung}]~
\end{description}
\defn{konfundierung}{\textbf{Konfundierung} liegt vor, wenn eine
Störvariable systematisch mit der unabhängigen Variablen variiert. Ihre
Wirkung auf die abhängige Variable ist dann von der Wirkung der UV nicht
mehr zu trennen.}
\bi
  \ite Beispiel: freiwillige App-Nutzer sind motivierter -- App oder
    Motivation?
  \ite die Täuschung in Variante~A: Längen und Messung lagen beim
    Vorführenden
\ei
\os{\newpage}

\begin{description}
\item[{\rot\bf Fünf Werkzeuge der Kontrolle}]~\\[-\lblskip]
\bi
  \ite \textbf{Randomisierung:} Zufall verteilt auch unbekannte
    Störvariablen
  \ite \textbf{Konstanthalten:} alles gleich bis auf die UV
  \ite \textbf{Schichten / Parallelisieren:} Gruppen nach bekanntem
    Störmerkmal ausbalancieren
  \ite \textbf{Kontrollgruppe}, möglichst \emph{aktiv}: gleich lang,
    gleich aufwendig, ohne Wirkstoff
  \ite \textbf{Blindbedingungen:} Teilnehmende und Messende kennen die
    Bedingung nicht
\ei
\os{\newpage}
\end{description}

%==============================================================================
\section{Externe Validität}
%==============================================================================
\nsl{Ein intern valides Ergebnis gilt zunächst nur für die untersuchten
Personen, unter den untersuchten Bedingungen, zum untersuchten Zeitpunkt.
\emph{Externe Validität} fragt, wie weit es sich auf andere Personen,
Situationen und Zeiten übertragen lässt \parencite{shadish2002experimental}.
Zwischen beiden besteht ein Zielkonflikt: Je strenger ein Labor alles
konstant hält, desto künstlicher wird die Situation. Das Buch nennt eine
besonders tückische Form dieses Problems: den \emph{Hot-Cold-Empathy-Gap}
\parencite{loewenstein1996visceral}. Wer im ruhigen Labor vorhersagen
soll, wie er unter Stress, Zeitdruck oder Erregung entscheiden würde,
beschreibt sein rationalisiertes Selbstbild. Umfragen zu zustandsabhängigem
Verhalten sind deshalb strukturell unzuverlässig, solange dieser Abstand
nicht methodisch adressiert wird (\buch{Kap.~4, HCEG}).}
\begin{description}
\item[{\rot\bf Externe Validität}]~\\[-\lblskip]
\bi
  \ite gilt das Ergebnis auch für andere Personen, Orte, Zeiten?
  \ite Zielkonflikt: je kontrollierter das Labor, desto künstlicher die
    Situation
  \ite Gelegenheitsstichproben (Freunde, Kommilitonen) sagen wenig über
    die Allgemeinbevölkerung
  \ite \textbf{Hot-Cold-Empathy-Gap}: im ruhigen Labor sagt man falsch
    voraus, wie man unter Stress handelt
\ei
\os{\newpage}
\end{description}

%==============================================================================
\section{Gütekriterien der Messung}
%==============================================================================
\nsl{Auch das beste Design hilft nicht, wenn die AV schlecht gemessen
wird. Die klassische Testtheorie nennt drei Hauptgütekriterien
\parencite{doering2016forschungsmethoden}. \emph{Objektivität}: Das
Ergebnis hängt nicht davon ab, wer die Messung durchführt, auswertet und
interpretiert. \emph{Reliabilität}: Die Messung ist zuverlässig, eine
Wiederholung unter gleichen Bedingungen liefert dasselbe Ergebnis.
\emph{Validität}: Die Messung erfasst tatsächlich das gemeinte Konstrukt.
Die Kriterien bauen aufeinander auf: Ohne Objektivität keine
Reliabilität, ohne Reliabilität keine Validität. Das Buch ergänzt diese
Begriffe in \buch{Kap.~10b} um ein praktisches Vokabular, das bei
schwer messbaren Konstrukten wie Staunen oder Stress hilft: Ein
\emph{Blind-Design} schaltet Erwartungen der Messenden aus. Ein
\emph{Prozessindikator} ist ein Relativmaß für Veränderung innerhalb einer
Person und einer Sitzung, kein Wahrheitsdetektor. Eine
\emph{Parallelmessung} erhebt dasselbe Konstrukt mit mehreren,
unabhängigen Verfahren (etwa Selbstbericht und Herzratenvariabilität),
damit ein Befund nicht an einem einzigen Instrument hängt. Auf der Ebene
des ganzen Designs heißt dieselbe Idee \emph{Triangulation}.}
\begin{description}
\item[{\rot\bf Drei Gütekriterien}]~\\[-\lblskip]
\begin{center}
\resizebox{0.86\textwidth}{!}{%
\begin{tikzpicture}[font=\sffamily,>=Stealth,
  st/.style={draw=rwuviolet,thick,fill=rwuviolet!8,rounded corners=3pt,
    minimum width=42mm,minimum height=17mm,align=center}]
  \node[st] (o) at (0,0) {\textbf{Objektivität}\\unabhängig davon,\\wer misst};
  \node[st] (r) at (5.4,0) {\textbf{Reliabilität}\\Wiederholung liefert\\dasselbe};
  \node[st,draw=rwucyan,fill=rwucyan40!40] (v) at (10.8,0) {\textbf{Validität}\\misst, was\\gemeint ist};
  \draw[->,very thick] (o) -- (r);
  \draw[->,very thick] (r) -- (v);
\end{tikzpicture}}
\end{center}
\bi
  \ite jedes Kriterium setzt das vorige voraus
  \ite Seile: Maßband (objektiv, reliabel) gegen Augenmaß des Vorführenden
\ei
\os{\newpage}

\item[{\rot\bf Messgüte-Vokabular aus dem Buch}]~\\[-\lblskip]
\bi
  \ite \textbf{Blind-Design:} wer misst, kennt die Bedingung nicht
  \ite \textbf{Prozessindikator:} Relativmaß für Veränderung in einer
    Person und Sitzung -- kein Wahrheitsdetektor
  \ite \textbf{Parallelmessung:} dasselbe Konstrukt mit mehreren
    Verfahren (Selbstbericht, Verhalten, Physiologie)
  \ite \textbf{Triangulation:} dieselbe Idee für das ganze Design --
    mehrere Methoden, Datenquellen, Perspektiven
\ei
\os{\newpage}
\end{description}

%==============================================================================
\section{Fallbeispiel: die Becher-Monty-Hall-Studie}
%==============================================================================
\nsl{Das Buch enthält ein vollständig ausgearbeitetes Forschungsdesign zu
seiner eigenen Leitthese, dass Staunen das Lernen begünstigt
(\buch{Kap.~12d}). Der Lerngegenstand ist das Monty-Hall-Problem, der
Staunmoment ein Becher-Zauberkunststück. Die Studie ist als
Zwischensubjekt-Design mit zufälliger Zuteilung angelegt. Beide Gruppen
sehen dasselbe, rund 13-minütige Erklärvideo; die Treatment-Gruppe sieht
zusätzlich, direkt nach dem Einstieg, einen rund zweiminütigen
Becher-Zaubertrick. Bei der Kontrollgruppe ist genau dieser Block
herausgeschnitten, alles Übrige ist identisch. Der Lehrteil nimmt mit
keinem Satz Bezug auf den Trick, damit er in beiden Fassungen wortgleich
und vollständig bleibt. Ein Manipulationscheck prüft, ob der Trick
überhaupt Staunen ausgelöst hat; gemessen werden außerdem Verständnis,
Transfer auf eine neu eingekleidete Aufgabe und Kalibrierung des eigenen
Urteils.

Der Entwurf benennt seinen Störfaktor selbst: Durch den reinen
Herausschnitt ist die Kontrollfassung rund zwei Minuten kürzer
(Treatment ca.~15:00, Kontrolle ca.~13:00). Länge und gemeinsame
Aufmerksamkeitszeit sind damit konfundiert. Die Abhilfe ist eine
längengleiche \emph{aktive Kontrolle}: An die Stelle des Tricks tritt ein
gleich langer, bewusst nicht staunenswerter Becher-Einschub, etwa ein
nüchternes Umstapeln mit sachlichem Kommentar
(\buch{Abschnitt~12d, Studienlogik}).}
\begin{description}
\item[{\rot\bf Der Idealentwurf (Kap.~12d)}]~\\[-\lblskip]
\bi
  \ite Zwischensubjekt-Design, zufällige Zuteilung
  \ite UV: Staunmoment (Becher-Zaubertrick) ja / nein
  \ite derselbe Film, ein Block herausgeschnitten -- sonst alles
    identisch
  \ite AV: Verständnis, Transfer, Kalibrierung
  \ite Manipulationscheck: Hat der Trick überhaupt Staunen ausgelöst?
\ei
\os{\newpage}

\item[{\rot\bf Der Störfaktor}]~\\[-\lblskip]
\begin{center}
\resizebox{0.9\textwidth}{!}{%
\begin{tikzpicture}[font=\sffamily]
  \def\u{0.6}
  \node[anchor=east] at (-0.2,0) {Treatment};
  \node[anchor=east] at (-0.2,-1) {Kontrolle};
  \node[anchor=east] at (-0.2,-2) {aktive Kontrolle};
  \fill[rwugray!40] (0,-0.3) rectangle (1.33*\u,0.3);
  \fill[rwuviolet] (1.33*\u,-0.3) rectangle (3.33*\u,0.3);
  \fill[rwugray!40] (3.33*\u,-0.3) rectangle (15*\u,0.3);
  \fill[rwugray!40] (0,-1.3) rectangle (13*\u,-0.7);
  \fill[rwugray!40] (0,-2.3) rectangle (1.33*\u,-1.7);
  \fill[rwucyan] (1.33*\u,-2.3) rectangle (3.33*\u,-1.7);
  \fill[rwugray!40] (3.33*\u,-2.3) rectangle (15*\u,-1.7);
  \node[anchor=west] at (15*\u+0.1,0) {ca.\ 15:00};
  \node[anchor=west] at (13*\u+0.1,-1) {ca.\ 13:00};
  \node[anchor=west] at (15*\u+0.1,-2) {ca.\ 15:00};
  \node[text=rwuviolet,above] at (2.33*\u,0.3) {Trick};
  \node[text=rwucyan,below] at (2.33*\u,-2.3) {neutraler Einschub};
\end{tikzpicture}}
\end{center}
\bi
  \ite Herausschneiden macht die Kontrolle \emph{kürzer}: Länge ist
    konfundiert
  \ite Abhilfe: gleich langer, nicht staunenswerter Einschub
\ei
\os{\newpage}
\end{description}

\nsl{Umgesetzt wurde der Entwurf als Online-Feldstudie mit drei Armen
(\buch{Kap.~12f}): U1 als Kontrolle ohne Zaubertrick, U2 mit dem
Becher-Zaubertrick als Spiel auf der Landingpage vor dem Fragebogen und
U3 mit dem in das Video integrierten, unaufgelösten Trick. Gegenüber dem
Idealentwurf ersetzt U2 die empfohlene längengleiche aktive Kontrolle
durch eine Platzierungsvariante. Die Zuteilung erfolgte zufällig über
eine Landingpage, die Feldzeit lief vom 3.~bis 31.~Juli 2026. Am Ende
lagen 32 auswertbare Fälle vor (13/9/10).

Die Befunde sind für eine Einheit über Forschungsdesign lehrreicher, als
es ein Erfolg gewesen wäre. Erstens ist der Manipulationscheck
gescheitert: Das Staunen liegt in den drei Armen bei 3{,}89 / 3{,}76 /
3{,}87, also praktisch gleichauf. Wenn die Intervention die vermutete
Ursache nicht messbar bewegt, sind alle nachgelagerten Fragen nach ihrer
Wirkung unbeantwortbar. Zweitens zeigte ein Zwischenstand am 16.~Juli
noch genau die erhoffte Abstufung (3{,}83 / 4{,}08 / 4{,}28), die sich
danach einebnete; hätte die Erhebung früher geendet, stünde ein schöner,
falscher Befund in der Welt. Drittens war der stärkste Zusammenhang einer,
den die Studie gar nicht untersuchen wollte: Wer das Monty-Hall-Problem
schon kannte, lag zu 92~Prozent richtig, wer es nicht kannte, zu rund der
Hälfte. Bei $n=32$ genügen wenige ungleich verteilte Vorwissensträger, um
jeden Armunterschied zu erzeugen oder zu tilgen. Viertens war die
Teststärke von Anfang an unzureichend: Für einen mittleren Effekt
($d=0{,}5$) hätte es rund 64 Personen pro Arm gebraucht, erreicht wurden
9 bis 13. Das Buch zieht daraus selbst die Konsequenzen für eine
Folgestudie: Fallzahl vorab planen, nach Vorwissen geschichtet
randomisieren und den Manipulationscheck reparieren.}
\begin{description}
\item[{\rot\bf Die Feldstudie (Kap.~12f)}]~\\[-\lblskip]
\bi
  \ite drei Arme: U1 Kontrolle, U2 Trick vorab, U3 Trick im Video
  \ite U2 statt der aktiven Kontrolle: eine Platzierungsvariante
  \ite 32 auswertbare Fälle (13 / 9 / 10)
  \ite Manipulationscheck gescheitert: Staunen 3{,}89 / 3{,}76 / 3{,}87
  \ite Zwischenstand 16.~Juli: 3{,}83 / 4{,}08 / 4{,}28 -- danach
    eingeebnet
\ei
\os{\newpage}

\item[{\rot\bf Was das Design daraus lernt}]~\\[-\lblskip]
\bi
  \ite \textbf{Vorwissen} überstrahlt alles: 92\,\% richtig mit, rund die
    Hälfte ohne Vorwissen
  \ite \textbf{Teststärke}: für $d=0{,}5$ rund 64 pro Arm nötig, erreicht
    9 bis 13
  \ite Folgestudie: Fallzahl vorab planen
  \ite nach Vorwissen \emph{geschichtet} randomisieren
  \ite Manipulationscheck reparieren (z.\,B.\ Verweildauer,
    verhaltensnahe Maße)
\ei
\os{\newpage}
\end{description}

\begin{merkbox}
Ohne Vergleich ist jede Beobachtung eine Vorführung. Ein Design ist so
gut wie die Alternativerklärungen, die es ausschließt -- und eine
Intervention ist erst dann geprüft, wenn der Manipulationscheck zeigt,
dass sie überhaupt gewirkt hat.
\end{merkbox}
\os{\newpage}

%==============================================================================
\section{Design-Werkstatt}
%==============================================================================
\begin{description}
\item[{\rot\bf Teil 1: Ein Design analysieren (Zweiergruppen, 15 Min.)}]~\\[-\lblskip]
\bi
  \ite Übungsblatt, Aufgabe~1: die (fiktive) Lern-App-Studie
  \ite Was sind UV und AV? Welcher Designtyp?
  \ite Gibt es eine Kontrollgruppe -- und taugt sie?
  \ite Welche Konfounder sehen Sie? Welche Bedrohungen der internen
    Validität?
  \ite Wie müsste das Design aussehen, damit die Schlussfolgerung trägt?
\ei
\os{\newpage}

\item[{\rot\bf Teil 2: Ihr eigenes Design (Einzelarbeit, 15 Min.)}]~\\[-\lblskip]
\bi
  \ite Ausgangspunkt: Ihre H\textsubscript{1} aus LE~5
  \ite Übungsblatt, Aufgabe~2: Design-Skizze ausfüllen
  \ite Messgüte-Check: objektiv? reliabel? valide? blind? parallel
    gemessen?
  \ite Welche Bedrohung der internen Validität bleibt -- und warum
    nehmen Sie sie in Kauf?
\ei
\os{\newpage}

\item[{\rot\bf Peer-Feedback (Zweiergruppen, 10 Min.)}]~\\[-\lblskip]
\bi
  \ite stellen Sie Ihr Design in drei Minuten vor
  \ite die Partnerin sucht \emph{eine} Alternativerklärung, die Ihr
    Design nicht ausschließt
  \ite gemeinsam: Welche Änderung würde sie ausschließen?
\ei
\os{\newpage}
\end{description}

\begin{lehrhinweis}
Bei der Analyse der fiktiven Studie finden die meisten Gruppen die
Selbstselektion schnell. Lohnend ist es, nach den weniger offensichtlichen
Konfoundern zu fragen: das Tutorium, das nur die App-Gruppe besucht hat,
und die nicht blinde Korrektur durch den App-Entwickler. Bei den eigenen
Designs ist der häufigste Fehler eine Kontrollgruppe, die sich in mehr als
einer Hinsicht von der Experimentalgruppe unterscheidet; hier hilft die
Frage \enquote{Was genau ist bei der Kontrollgruppe anders, und ist das
\emph{nur} Ihre UV?}. Wer ein nicht-experimentelles Design wählt, muss das
nicht ändern, sollte aber in der Begründung zeigen, dass ihm die Grenzen
der Kausalaussage klar sind. Statt der fiktiven Studie kann auch der
Methodenteil eines echten Artikels analysiert werden; als realer Fall aus
dem Buch eignet sich die Feldstudie aus \buch{Kap.~12f}, die ihre
Abweichungen vom Idealentwurf selbst offenlegt.
\end{lehrhinweis}

%==============================================================================
\section{Lernziele und Aufgaben}
%==============================================================================
\begin{description}
\item[{\rot\bf Nach dieser Einheit können Sie}]~\\[-\lblskip]
\bi
  \ite UV, AV und Störvariablen in einem Design benennen
  \ite experimentelle, quasi-experimentelle und nicht-experimentelle
    Designs unterscheiden
  \ite Bedrohungen der internen Validität erkennen und Gegenmittel
    benennen
  \ite Objektivität, Reliabilität und Validität einer Messung prüfen
  \ite ein angemessenes Design für Ihre Fragestellung auswählen und
    begründen
\ei
\os{\newpage}

\item[{\rot\bf Bis zur nächsten Woche}]~\\[-\lblskip]
\bi
  \ite \textbf{Fallstudie} (Aufgabe~1) fertigstellen
  \ite \textbf{Design-Skizze} (Aufgabe~2): Design für Ihre eigene
    Hypothese, Begründung auf höchstens einer Seite -- benotete
    Teilleistung
  \ite Feedback Ihrer Partnerin einarbeiten
  \ite \textbf{Tagebuch des Staunens:} eine Behauptung aus Werbung oder
    Nachrichten, der ein Vergleich fehlt
\ei
\os{\newpage}
\end{description}

\handoutinput{handouts/LE06}
